\documentclass[11pt,a4paper]{article}
\usepackage[utf8]{inputenc}
\usepackage[german]{babel}
\usepackage{amsmath, amssymb, amsthm}
\usepackage{geometry}
\geometry{margin=2.5cm}

\title{\textbf{Binäre Resonanz im Pascal-Dreieck: Das Produkt von Primzahlvierlingen als harmonischer Grundbaustein}}
\author{Thomas Hoffbauer}
\date{April 2026}

\begin{document}
	
	\maketitle
	
	\begin{abstract}
		Diese Arbeit untersucht die numerische Resonanz zwischen den Produkten von Primzahlvierlingen und der fraktalen Binärstruktur des Pascalschen Dreiecks. Wir zeigen, dass das Produkt des kleinsten Vierlings $P_1 = 5 \cdot 7 \cdot 11 \cdot 13 = 5005$ eine fundamentale Symmetrie zur 15. Schicht des Dreiecks aufweist. Mittels des Satzes von Legendre und einer Analyse des Hamming-Gewichts wird ein Modell vorgeschlagen, das Primzahlcluster als harmonische Oberschwingungen einer binären Grundfrequenz interpretiert.
	\end{abstract}
	
	\section{Motivation und Einleitung}
	Die Verteilung von Primzahlclustern, insbesondere von Primzahlvierlingen der Form $(p, p+2, p+6, p+8)$, gilt klassischerweise als stochastisch. Die vorliegende Untersuchung bricht mit dieser Sichtweise und postuliert eine deterministische Ordnung innerhalb der kombinatorischen Gitterstruktur des Pascalschen Dreiecks (PD). 
	
	Die Kernidee ist, dass Primzahlprodukte keine isolierten Werte sind, sondern als konstruktive Interferenzmuster in einer binären Matrix agieren. Die Zahl $5005$ dient hierbei als „genetischer Code“ oder Grundbaustein für die Lokalisierung höherer Cluster.
	
	\section{Der Grundbaustein: Die 5005-Resonanz}
	Wir betrachten das Produkt $P_1$:
	$$P_1 = 5 \cdot 7 \cdot 11 \cdot 13 = 5005$$
	Im Pascalschen Dreieck entspricht dieser Wert dem Binomialkoeffizienten:
	$$\binom{15}{6} = 5005$$
	Die Schicht $n=15$ ist im Binärsystem als $1111_2$ darstellbar. Dies ist ein Zustand maximaler Sättigung, in dem alle Binomialkoeffizienten ungerade sind. 
	
	\section{Mathematisches Framework}
	
	\subsection{Die erzeugende Exponentialfunktion}
	Als Ausgangspunkt nutzen wir die Identität:
	$$(e^x - 1)^n = \sum_{i=0}^{n} \binom{n}{i} e^{ix} (-1)^{n-i}$$
	Wir interpretieren $x$ als Trägerwelle der Primzahlstruktur und $n$ als binären Filter. Eine Resonanz tritt auf, wenn die Wechselsumme über die Binomialkoeffizienten eine minimale binäre Entropie aufweist.
	
	\subsection{Der Satz von Legendre und die binäre Dichte}
	Um die Position im Sierpinski-Gitter analytisch zu fassen, nutzen wir den Satz von Legendre für die $p$-adische Bewertung von Fakultäten bei $p=2$:
	$$v_2(n!) = n - w(n)$$
	wobei $w(n)$ das Hamming-Gewicht (Anzahl der Einsen) von $n$ ist. Für den Binomialkoeffizienten ergibt sich die „binäre Spannung“:
	$$v_2\left(\binom{n}{k}\right) = w(k) + w(n-k) - w(n)$$
	Für den Grundbaustein $P_1 = \binom{15}{6}$ gilt $v_2 = w(6) + w(9) - w(15) = 2 + 2 - 4 = 0$. Dies beweist, dass $5005$ eine ungerade Zahl an einem Ort maximaler fraktaler Dichte ist.
	
	\section{Differenzgleichung und Rhythmus}
	Wir postulieren, dass jeder weitere Primzahlvierling $P_m$ eine Resonanzbedingung erfüllt, die durch die Verschiebung der Bit-Muster definiert ist:
	$$\Delta \Omega = | w(P_m) - \Phi(n) | \approx 0$$
	Die Abfolge der Vierlinge folgt einer harmonischen Reihe im Binärraum, wobei die Frequenz durch die Autokorrelation der Bit-Sequenzen bestimmt wird:
	$$\rho(\tau) = \sum B_i \cdot B_{i+\tau}$$
	
	\section{Schlussfolgerung}
	Die Analyse zeigt, dass Primzahlvierlinge als „digitale Echos“ des Grundbausteins 5005 verstanden werden können. Die Transformation von der klassischen Zahlentheorie hin zu einer Signalverarbeitung auf dem Pascal-Dreieck ermöglicht es, Cluster-Positionen durch binäre Resonanzfilterung vorherzusagen. Die 5005 ist in diesem System die harmonische Grundwelle, auf der die gesamte Architektur der Vierlinge aufbaut.
	
\end{document}